\section{Анализ предметной области разрабатываемой системы}

\subsection{Особенности предметной области и назначение разрабатываемого web-сервиса}

В данной работе рассматривается разработка web-сервиса, предназначенного для размещения услуг, поиска исполнителей и организации записи на свободное время. Разрабатываемая система ориентирована на две основные категории пользователей: исполнителей, которые публикуют свои услуги и управляют расписанием, и заказчиков, заинтересованных в поиске подходящего предложения и бронировании удобного временного интервала.

Сфера оказания бытовых, цифровых и частных услуг в настоящее время развивается особенно активно. Пользователи всё чаще выбирают не традиционный способ поиска специалиста через знакомых, объявления или переписку в мессенджерах, а цифровые платформы, позволяющие в короткий срок найти подходящего исполнителя, сравнить предложения и выбрать наиболее удобный вариант. При этом заказчик ожидает от сервиса не только наличия каталога услуг, но и понятного механизма выбора даты, времени, стоимости и условий взаимодействия.

Для исполнителя наличие специализированного web-сервиса также имеет существенное значение. Ему необходимо не только размещать информацию об услуге, но и управлять своей занятостью, поддерживать порядок в записях, быстро реагировать на новые заявки и контролировать уже подтверждённые бронирования. При отсутствии автоматизации эти процессы часто выполняются вручную, что приводит к избыточным трудозатратам и увеличивает вероятность организационных ошибок.

Разрабатываемый сервис объединяет основные действия в рамках единой информационной среды. Исполнитель получает возможность создать карточку услуги, указать категорию, стоимость, продолжительность, описание и изображения, а также задать рабочее расписание на несколько недель вперёд. Заказчик, в свою очередь, может просматривать каталог услуг, применять фильтрацию по категориям, городу и датам, выбирать свободные слоты и оформлять бронирование без длительного предварительного согласования.

Специфика данной предметной области заключается в том, что ключевую роль здесь играет управление временем. В отличие от обычной доски объявлений, сервис размещения услуг должен учитывать не только наличие самой услуги, но и фактическую доступность исполнителя в конкретные дни и часы. Если пользователь видит предложение, но не может выбрать подходящий слот, эффективность платформы заметно снижается. По этой причине важнейшими элементами системы становятся расписание, слоты записи, статусы заявок и быстрый доступ к информации о доступных интервалах.

Дополнительной особенностью рассматриваемой предметной области является наличие взаимной ответственности сторон. При работе с записью на услуги большое значение имеют не только отзывы об исполнителе, но и сведения о добросовестности заказчика. В практической деятельности возможны отмены записи, переносы, неявки клиентов и иные ситуации, которые влияют на занятость исполнителя и качество организации работы. Следовательно, система должна учитывать не только сам факт бронирования, но и результаты взаимодействия между пользователями.

Таким образом, предметная область включает процессы публикации услуг, ведения расписания, поиска по параметрам, бронирования свободных интервалов, переписки по заявкам и накопления репутационной информации. Это определяет необходимость разработки специализированного web-сервиса, ориентированного на комплексную организацию записи и сопровождение оказания услуг.

\subsection{Анализ существующих проблем при ручной организации записи на услуги}

Традиционный подход к организации записи на услуги во многих случаях остаётся разрозненным и недостаточно удобным. Исполнитель публикует информацию о себе на одной площадке, принимает обращения через телефон или мессенджеры, ведёт расписание вручную в заметках, таблицах или блокноте, а сведения о договорённостях и клиентах хранит в личной переписке. Такой способ может быть приемлем на раннем этапе, однако при увеличении числа заказов начинает создавать заметные трудности.

Одной из главных проблем является отсутствие единого пространства хранения информации. Данные о клиентах, стоимости, времени записи, пожеланиях по заказу и статусе обращения находятся в разных источниках. Это затрудняет поиск нужной информации, создаёт риск потери части сведений и усложняет повторное взаимодействие с заказчиком. Исполнителю приходится тратить время не на выполнение услуги, а на восстановление организационной информации.

Серьёзной проблемой является и ручное ведение расписания. При таком подходе исполнитель самостоятельно отслеживает свободные интервалы, согласовывает их с клиентами и фиксирует изменения без автоматической проверки пересечений. В результате возможны ошибки при записи: наложение нескольких заявок на одно время, потеря свободных окон, неверный отказ клиенту или случайное дублирование бронирований. Чем больше услуг и клиентов, тем выше вероятность подобных ошибок.

Отдельное затруднение связано с тем, что заказчик часто не видит реальную доступность исполнителя. На большинстве площадок пользователь может только просмотреть услугу и написать исполнителю. Далее начинается переписка, в которой стороны поэтапно согласовывают дату, время и условия. Такой процесс занимает дополнительное время, снижает удобство использования сервиса и в ряде случаев приводит к тому, что заказчик отказывается от дальнейшего взаимодействия.

Проблемой также является отсутствие структурированного сопровождения уже созданной записи. После предварительного согласования необходимо хранить информацию о статусе заявки, времени бронирования, комментариях клиента, изменениях условий и результате оказания услуги. Если все эти данные остаются только в сообщениях, их трудно использовать повторно, а сама коммуникация быстро становится перегруженной и малоуправляемой.

Ещё одной важной проблемой является недостаток репутационной информации о заказчике. Во многих сервисах основное внимание уделяется только рейтингу исполнителя, тогда как у исполнителя отсутствует удобный механизм фиксации неявки клиента, срыва записи или недобросовестного поведения. Это снижает прозрачность работы платформы и не позволяет учитывать предыдущий опыт взаимодействия при обработке новых заявок.

Кроме того, ручной подход плохо масштабируется. При небольшом количестве заказов исполнитель ещё может контролировать ситуацию самостоятельно, однако при росте клиентской базы, увеличении числа услуг и расширении расписания затраты времени заметно возрастают. Таким образом, традиционный способ организации записи не соответствует современным требованиям к удобству, скорости и надёжности взаимодействия между сторонами.

\subsection{Причины необходимости разработки автоматизированного решения}

Указанные недостатки ручной организации записи могут быть устранены путём разработки специализированного web-сервиса, автоматизирующего ключевые процессы взаимодействия между исполнителями и заказчиками.

Прежде всего автоматизация должна обеспечить централизованное хранение информации о пользователях, услугах, категориях, городах, расписании, заявках, сообщениях и отзывах. Это позволит отказаться от использования нескольких независимых инструментов и организовать согласованную работу с данными в рамках одной системы.

Особое значение имеет автоматизация календарного планирования. Исполнитель должен иметь возможность задавать интервалы работы по дням недели, формировать расписание на несколько недель вперёд и отмечать недоступные даты. На основе этих данных система должна автоматически создавать свободные слоты и предоставлять заказчику только актуальные варианты для записи.

Не менее важной задачей является автоматизация поиска услуг по доступным датам. Заказчик должен иметь возможность выбрать не одну жёстко заданную дату, а несколько подходящих дней, после чего система должна показать только те услуги, по которым в выбранные даты действительно существует свободное время. Такой подход повышает удобство поиска и позволяет быстрее принимать решение.

Автоматизированное решение должно также обеспечивать полноценное бронирование услуги. После выбора слота система должна сохранять дату, время и статус заявки, связывать заказчика и исполнителя в рамках одного бронирования и поддерживать дальнейшее изменение состояний, например подтверждение, отмену, завершение услуги или фиксацию неявки.

Для удобства последующего взаимодействия необходим встроенный чат, связанный с конкретным бронированием. Такой механизм позволяет хранить переписку структурированно, не смешивая сообщения по разным услугам и заявкам, а также повышает прозрачность общения по каждому отдельному заказу.

Дополнительным аргументом в пользу автоматизации является учёт репутационной информации. В системе должен быть реализован механизм оценки добропорядочности заказчика на основе отзывов исполнителей после завершённой услуги либо неявки клиента. Это позволит учитывать историю предыдущих взаимодействий и повысит надёжность работы платформы.

Следовательно, разработка автоматизированного решения необходима для повышения удобства поиска и записи на услуги, снижения числа организационных ошибок, улучшения управления расписанием исполнителя, структурирования коммуникации между сторонами и повышения прозрачности всего процесса оказания услуг.

\subsection{Сравнительный анализ существующих интернет-сервисов в рассматриваемой области}

В настоящее время существует ряд цифровых платформ, которые частично или полностью решают задачи поиска исполнителей, организации записи и сопровождения взаимодействия между сторонами. Однако каждая из них ориентирована на собственную модель работы и не всегда соответствует требованиям, которые предъявляются к разрабатываемому web-сервису. Для определения особенностей будущей системы целесообразно рассмотреть наиболее близкие по назначению решения.

\subsubsection{Анализ возможностей сервиса YCLIENTS}

YCLIENTS представляет собой платформу онлайн-записи и автоматизации работы сервисного бизнеса. Основной акцент в данном решении сделан на управление расписанием, онлайн-запись клиентов, напоминания, уведомления, предоплату и работу с клиентской базой. Такой подход удобен для компаний и специалистов, деятельность которых строится вокруг приёма клиентов по расписанию.

Преимуществом данного сервиса является развитый инструментарий календарного управления. Система позволяет организовать запись на конкретные услуги, поддерживать расписание сотрудников, использовать онлайн-платежи, напоминания и личный кабинет клиента. Для организаций, которым требуется автоматизация уже сформированной клиентской записи, такой функционал является весьма эффективным.

Однако YCLIENTS ориентирован преимущественно на автоматизацию работы отдельной компании, салона или студии, а не на универсальную площадку взаимодействия множества независимых исполнителей и заказчиков. В рамках рассматриваемой дипломной разработки требуется не только запись в уже существующий сервисный бизнес, но и публикация услуг разными исполнителями в общем каталоге, поиск по доступным датам и выбор специалиста пользователем в рамках одной платформы. Следовательно, функциональность YCLIENTS полезна как ориентир в части календарной записи, но не полностью соответствует модели маркетплейса услуг.

\subsubsection{Анализ возможностей платформы Профи.ру}

Профи.ру представляет собой сервис, основной задачей которого является организация встречи клиента и специалиста. Пользователь формирует запрос на нужную услугу, после чего получает возможность подобрать подходящего исполнителя, ориентируясь на профиль, отзывы, опыт и иные параметры.

Преимуществом данного сервиса является развитый механизм подбора специалистов по большому числу категорий услуг. Платформа позволяет быстро находить исполнителей в различных областях, сравнивать предложения и принимать решение на основе репутационных характеристик. Для заказчика такой подход удобен в ситуациях, когда необходимо выбрать специалиста по нескольким критериям одновременно.

В то же время данная модель в большей степени ориентирована на поиск и выбор исполнителя, чем на непосредственную запись на заранее опубликованную услугу с жёстко заданным расписанием. В рамках разрабатываемой системы важно не только найти подходящего специалиста, но и сразу выбрать конкретную дату и время оказания услуги. Следовательно, Профи.ру является полезным примером сервиса поиска исполнителей, однако не закрывает в полном объёме задачу календарного бронирования слотов.

\subsubsection{Анализ возможностей сервиса YouDo}

YouDo представляет собой сервис, основанный на размещении задания заказчиком и получении откликов от исполнителей. Пользователь описывает задачу, после чего сервис позволяет выбрать исполнителя среди откликнувшихся кандидатов с учётом отзывов, цены и других характеристик.

Сильной стороной данного решения является конкурентная модель выбора исполнителя. Заказчик получает несколько вариантов, может сравнить отклики и выбрать оптимальное предложение. Сервис также ориентирован на широкий спектр бытовых и деловых услуг, что делает его показательным примером платформы для поиска специалистов.

Однако логика работы YouDo строится прежде всего вокруг задания и откликов, а не вокруг публикации готовой карточки услуги с календарём доступности. Для разрабатываемого web-сервиса требуется иной сценарий взаимодействия: исполнитель заранее публикует услугу, указывает параметры предложения, формирует расписание, а заказчик выбирает не абстрактного исполнителя под задачу, а конкретную услугу и сразу оформляет запись на свободный слот. По этой причине YouDo не может быть использован как прямой аналог разрабатываемой системы, хотя и представляет интерес с точки зрения организации рыночного взаимодействия между сторонами.

\subsubsection{Итоги анализа существующих решений}

Проведённый анализ показал, что существующие платформы решают отдельные задачи рассматриваемой предметной области, но не обеспечивают в полном объёме ту совокупность функций, которая требуется для разрабатываемого web-сервиса.

Сервисы типа YCLIENTS эффективно автоматизируют расписание, запись клиентов и внутренние процессы конкретного бизнеса, однако не ориентированы на универсальную площадку размещения услуг от множества независимых исполнителей.

Платформы типа Профи.ру и YouDo, напротив, хорошо решают задачу поиска специалиста, выбора исполнителя и учёта отзывов, но в меньшей степени ориентированы на календарную запись с отображением реально доступных временных интервалов и дальнейшим сопровождением бронирования в едином интерфейсе.

Таким образом, для рассматриваемой предметной области целесообразной является разработка собственного специализированного решения, объединяющего публикацию услуг, поиск по параметрам, фильтрацию по доступным датам, бронирование свободных слотов, встроенный чат и механизм оценки добропорядочности заказчика. Именно такое объединение функций позволит получить систему, максимально соответствующую требованиям проекта.

\subsection{Определение основных задач разработки}

На основании проведённого исследования предметной области можно сформулировать основные задачи, которые должна решать разрабатываемая система:

\begin{enumerate}
	\item Обеспечить регистрацию и авторизацию пользователей с разграничением ролей заказчика и исполнителя.
	\item Реализовать механизм создания, хранения и редактирования карточек услуг с указанием категории, описания, стоимости, длительности и изображений.
	\item Реализовать механизм ведения расписания исполнителя с возможностью задания рабочих интервалов по дням недели на несколько недель вперёд.
	\item Обеспечить автоматическое формирование и отображение свободных слотов для записи.
	\item Реализовать поиск услуг по категориям, городу и доступным датам с отображением только тех предложений, по которым имеется свободное время.
	\item Реализовать создание заявок на бронирование услуги с сохранением даты, времени и статуса записи.
	\item Обеспечить изменение состояний заявки, включая подтверждение, отмену, завершение услуги и фиксацию неявки заказчика.
	\item Реализовать встроенный чат в рамках конкретного бронирования для взаимодействия между исполнителем и заказчиком.
	\item Обеспечить механизм добавления услуг в избранное для последующего быстрого доступа.
	\item Реализовать систему отзывов и механизм оценки добропорядочности заказчика по результатам оказания услуги.
	\item Обеспечить хранение данных в базе данных и предоставить удобный web-интерфейс для работы с системой.
\end{enumerate}